\chapter{Stochastic Spiking and Neural Noise}
\label{ch:stochastic-spiking}

\textit{Deterministic membrane equations predict a unique trajectory for each input, but real neurons fire variably even under repeated stimulation. This chapter introduces the statistical language needed to analyze stochastic spike trains: interspike-interval distributions, count variability, spectral measures, and signal-transmission metrics such as coherence and information rate. The goal is to connect single-neuron noise statistics to computational function.}

% ============================================================
\section{Why a Stochastic Description?}
\label{sec:why-stochastic}
% ============================================================

In \cref{ch:lif,ch:hodgkin-huxley,ch:excitability}, spikes emerged from deterministic dynamics. If we repeat the same current pulse in those models, we get the same spike train each trial. Real neurons are different: repeated trials show jitter in spike timing, variable spike counts, and spontaneous firing without explicit stimulus.

The sources of variability are multiple: channel noise from stochastic gating (\cref{ch:single-channels}), synaptic release noise, fluctuating network input, and hidden state variables not included in reduced models. A deterministic equation remains useful for the mean trajectory, but statistical descriptors are needed for variability around that mean.

Computationally, this shift matters because information is carried not only by mean firing rates but also by timing precision and trial-to-trial variability. Two neurons with identical mean rates can have very different coding performance if their spike statistics differ.

% ============================================================
\section{Spike Trains as Point Processes}
\label{sec:point-process}
% ============================================================

A spike train is represented as a sum of Dirac impulses at spike times $\{t_k\}$:
%
\begin{equation}
  x(t) = \sum_k \delta(t - t_k)
  \label{eq:spike-train-delta}
\end{equation}
%
This turns event times into a signal suitable for averaging, spectrum analysis, and linear-response theory.

Define the interspike intervals (ISIs) $T_k = t_{k+1}-t_k$. Their distribution $p_{\text{ISI}}(T)$ is the first statistical object. The mean firing rate is
%
\begin{equation}
  r = \frac{1}{\langle T \rangle}.
  \label{eq:rate-isi}
\end{equation}
%

The coefficient of variation of ISIs is
%
\begin{equation}
  \mathrm{CV} = \frac{\sqrt{\Var[T]}}{\langle T \rangle}.
  \label{eq:cv}
\end{equation}
%
For near-periodic firing, $\mathrm{CV}\ll 1$; for irregular firing, $\mathrm{CV}\sim 1$; for bursty or heavy-tailed statistics, $\mathrm{CV}>1$.

Computationally, ISI statistics separate distinct coding regimes: reliable temporal coding requires low-jitter ISIs, while high-variability regimes can support probabilistic sampling or exploration.

% ============================================================
\section{Count Variability and the Fano Factor}
\label{sec:fano}
% ============================================================

Let $N_T$ be the number of spikes in a counting window of duration $T$. The Fano factor is
%
\begin{keyequation}[Fano Factor]
\begin{equation}
  F(T) = \frac{\Var[N_T]}{\mathbb{E}[N_T]}.
  \label{eq:fano}
\end{equation}
\tcblower
$F=1$: Poisson-like variability, \quad
$F<1$: more regular than Poisson, \quad
$F>1$: overdispersed (e.g., bursts, slow-rate fluctuations).
\end{keyequation}

For renewal processes, short and long windows can yield different behavior; slow gain fluctuations typically inflate $F(T)$ at large $T$. This scale dependence is important experimentally: a neuron may look near-Poisson at 20~ms and strongly overdispersed at 1~s.

Computationally, the Fano factor measures reliability of rate codes. Lower $F$ improves detectability of weak stimulus-driven rate changes in finite observation windows.

% ============================================================
\section{Spectra, Coherence, and Signal Transmission}
\label{sec:spectra-coherence}
% ============================================================

When spikes encode time-varying signals, second-order frequency-domain statistics become central. Let $s(t)$ be a stimulus and $x(t)$ the spike train. Define auto- and cross-spectra $S_{xx}(f)$, $S_{ss}(f)$, and $S_{xs}(f)$.

The coherence between stimulus and spike train is
%
\begin{equation}
  C_{xs}(f) = \frac{|S_{xs}(f)|^2}{S_{xx}(f)S_{ss}(f)},
  \qquad 0 \le C_{xs}(f) \le 1.
  \label{eq:coherence}
\end{equation}
%

$C_{xs}(f)$ quantifies how well spike fluctuations at frequency $f$ are locked to stimulus fluctuations at the same frequency. The phase of $S_{xs}(f)$ gives the frequency-dependent lag.

A commonly used lower bound on mutual-information rate for Gaussian channels is
%
\begin{equation}
  I \ge -\int_0^{\infty} \log_2\!\bigl(1-C_{xs}(f)\bigr)\,\dd f.
  \label{eq:info-coherence}
\end{equation}
%

Computationally, this framework makes temporal coding measurable: resonator neurons often show coherence peaks near preferred frequencies, while integrator-like neurons transmit more broadband low-frequency content.

% ============================================================
\section{Noise-Driven Spiking in Reduced Models}
\label{sec:noise-driven-lif}
% ============================================================

A minimal stochastic extension of the LIF model adds a fluctuating input term:
%
\begin{equation}
  \taum \frac{\dd V}{\dd t} = -(V-\Erev{L}) + R_m\bigl[I_0 + \sigma\,\xi(t)\bigr],
  \label{eq:stoch-lif}
\end{equation}
%
with threshold/reset as usual and white noise $\xi(t)$ satisfying
$\langle \xi(t) \rangle=0$ and
$\langle \xi(t)\xi(t') \rangle=\delta(t-t')$.

Even when $I_0$ is subthreshold, noise can trigger spikes (escape over an effective barrier). The result is a nonzero spontaneous firing rate with broad ISI statistics. As $\sigma$ increases, CV and Fano factor typically increase, and timing precision relative to fast signals decreases.

Computationally, this connects dynamical excitability to variability statistics: bifurcation type sets deterministic backbone, while noise sets operating-point reliability around that backbone.

% ============================================================
\section{Spontaneous vs Evoked Activity}
\label{sec:spont-vs-evoked}
% ============================================================

In spontaneous activity, statistics are characterized without controlled input: baseline rate, ISI density, CV, Fano factor, and spike-train spectrum. In evoked activity, we additionally measure transfer properties: gain, phase lag, coherence, and information rate.

A practical workflow is:
(1) estimate baseline renewal statistics,
(2) apply weak sinusoidal or stochastic stimuli,
(3) compute linear response and coherence,
(4) compare observed coding limits to predictions from the baseline noise model.

This workflow links what looks like ``noise'' to what is actually an important determinant of neural computation. Noise sets detection thresholds, response latency distributions, and the precision with which time-varying stimuli can be represented.


\begin{chaptersummary}

\begin{center}
\renewcommand{\arraystretch}{1.6}
\begin{tabularx}{\textwidth}{@{}l X X@{}}
  \toprule
  \textbf{Concept} & \textbf{Equation} & \textbf{Key Insight} \\
  \midrule
  Spike train representation &
  $x(t)=\sum_k \delta(t-t_k)$ &
  Event times as analyzable signal \\
  %
  Rate from ISIs &
  $r = 1/\langle T\rangle$ &
  Mean interval sets mean firing rate \\
  %
  ISI irregularity &
  $\mathrm{CV}=\sqrt{\Var[T]}/\langle T\rangle$ &
  Dimensionless timing variability \\
  %
  Count variability &
  $F(T)=\Var[N_T]/\mathbb{E}[N_T]$ &
  Reliability of rate code in finite windows \\
  %
  Coherence &
  $C_{xs}(f)=|S_{xs}|^2/(S_{xx}S_{ss})$ &
  Frequency-specific stimulus locking \\
  %
  Stochastic LIF &
  $\taum\dot V=-(V-\Erev{L})+R_m(I_0+\sigma\xi)$ &
  Noise converts deterministic threshold to probabilistic spiking \\
  \bottomrule
\end{tabularx}
\end{center}

\end{chaptersummary}

\begin{conceptquestions}
  \item Two neurons have the same mean firing rate but different CV values ($0.2$ vs $1.1$). What distinct coding regimes do these values suggest?

  \item Why can the Fano factor depend strongly on counting-window duration $T$? What mechanism produces large-$T$ overdispersion?

  \item Coherence $C_{xs}(f)$ is high near 10~Hz and low elsewhere. What does this say about the neuron's temporal selectivity and likely excitability class?

  \item In the stochastic LIF model, how can a neuron fire when $I_0$ is subthreshold? How does increasing $\sigma$ qualitatively affect ISI statistics?

  \item Why is it useful to analyze both spontaneous and evoked activity rather than only one of them?
\end{conceptquestions}
